\section{Customer-Owned Model Hosting and Compute Deployment}
\label{sec:aion-model-hosting-compute-deployment}

\subsection{Purpose and ownership principle}

The user owns the AION brain and chooses where the models providing its
computational intelligence run.  AION is not inseparable from any particular
language model, cloud provider, hardware vendor, or inference engine.  The
customer's identity, memory, business map, policy, permissions, learned
procedures, and outcome history remain under customer control even when a
remote model supplies additional computation.

The central ownership rule is:

\begin{quote}
The AION brain is the persistent intelligence and authority system.  A model is
replaceable cognitive compute used by that brain for a bounded task.
\end{quote}

This separation allows a household or organisation to begin with a modest local
model, add private cloud capacity when needed, use an external premium provider
for selected work, and replace any of those components without losing what AION
has learned.

\subsection{Primary deployment configurations}

\subsubsection{Everything local}

In the fully local configuration, the AION brain and a suitably sized
open-weight model run on the customer's personal computer, workstation, or home
or business server.  The operating path is:

\[
\text{Pilot}
\rightarrow \text{AION brain}
\rightarrow \text{local model}
\rightarrow \text{AION verification}
\rightarrow \text{Pilot or governed action}.
\]

This configuration requires no model API and can keep all working data on the
customer's machine.  After the hardware has been acquired, the marginal cost of
ordinary inference may be low, although electricity, maintenance, and hardware
depreciation still exist.  Capacity is limited by processor speed, accelerator
availability, memory, model size, and the number of concurrent users.

The fully local configuration is appropriate for households, sole traders,
small teams, sensitive offline work, and environments in which internet access
is unavailable or prohibited.  AION must remain genuinely useful in this mode;
cloud access may improve quality or speed but must not be the sole source of
core intelligence.

\subsubsection{Customer-owned cloud compute}

In the customer-owned cloud configuration, AION may remain on a local mother
brain or run inside infrastructure controlled by the customer's organisation.
The customer deploys an approved model to its own account with a cloud,
infrastructure, or confidential-compute provider.  Examples include a private
AWS, Google Cloud, or Azure environment, NVIDIA infrastructure, Prem, or another
qualified provider.

The logical path is:

\[
\begin{split}
\text{Pilot}
&\rightarrow \text{customer's AION brain} \\
&\rightarrow \text{minimum permitted task context} \\
&\rightarrow \text{customer's private model endpoint} \\
&\rightarrow \text{structured result} \\
&\rightarrow \text{AION validation, policy, and execution}.
\end{split}
\]

The model may run on a rented GPU instance, a managed container, a private
inference cluster, or customer-owned hardware.  Deployment may use an NVIDIA
NIM container, a Tessaris-qualified model package, or another compatible
inference server.  AION stores the endpoint configuration and credentials in
the customer's private vault and calls the endpoint only when policy permits
the disclosure and the task requires the additional capacity.

The remote model returns a proposed answer, classification, plan, extraction,
comparison, or tool instruction.  It does not inherit authority over the
business and cannot directly rewrite canonical memory, grant permissions, or
perform consequential actions merely because it produced a response.

For stronger isolation, the AION brain, business data, retrieval services, and
model endpoint may all run inside the same customer-controlled virtual private
cloud.  In that arrangement, authorised business context can remain inside the
customer's network boundary.  Customer-owned cloud deployment does not by
itself guarantee privacy: identity controls, network isolation, encryption,
logging, retention, administrator access, backups, and regional processing must
also be configured and verified.

\subsubsection{External model API}

In the external-provider configuration, the user supplies credentials for an
approved service such as Gemini, OpenAI, Anthropic, or another compatible
provider.  The operating path is:

\[
\text{Pilot}
\rightarrow \text{AION disclosure policy}
\rightarrow \text{external model API}
\rightarrow \text{result}
\rightarrow \text{AION verification}.
\]

This option is convenient because the customer need not administer GPU
infrastructure.  However, the permitted request context is transmitted to the
selected third party.  Before transmission, AION must apply the customer's
data-classification, purpose, residency, retention, provider, and cost policies.
Where the disclosure is material, the user must be shown what category of data
will leave the customer boundary, which provider will receive it, and why that
provider was selected.

External APIs are optional quality or capacity providers rather than roots of
AION identity or memory.  Provider failure, quota exhaustion, key removal, or a
commercial change must cause a policy-compliant fallback or an honest
limitation; it must not corrupt the customer brain.

\subsection{What the customer downloads}

The initial Tessaris package should contain or install the following components:

\begin{itemize}
  \item the AION brain and governed learning runtime;
  \item the Personal Pilot, Workspace, and Boardroom interfaces permitted by
        the selected edition;
  \item identity, possession, permissions, delegation, and approval controls;
  \item encrypted memory and the permission-aware business map;
  \item the model, evidence, capability, and compute router;
  \item a small local model appropriate for the detected hardware, or a guided
        installer for that model;
  \item connectors for optional private compute and external providers;
  \item capability packages for files, communications, business systems, and
        authorised devices;
  \item health monitoring, signed updates, automatic recovery, and rollback;
  \item encrypted backup, verified restore, complete export, and migration
        controls.
\end{itemize}

The customer does not necessarily download a very large model to a laptop.
During setup, AION profiles the available processor, memory, accelerator,
storage, power constraints, and expected workload.  It then explains the
appropriate options in non-technical language.  A representative prompt is:

\begin{quote}
This computer can comfortably run a small local model for everyday private
work.  Larger business analysis will be faster with additional private compute.
Would you like to connect infrastructure in an account you control?
\end{quote}

If the customer selects AWS, Google Cloud, Azure, NVIDIA, Prem, or another
qualified target, a guided deployment creates or identifies the model endpoint
within that account and registers it with AION.  Provider credentials remain in
the customer-controlled vault.  The television, phone, and ordinary device
nodes never receive raw cloud secrets.

\subsection{Component and authority separation}

The architecture separates durable customer state from replaceable compute:

\begin{description}
  \item[AION brain.] Owns memory, the business map, policy, authority, learning,
        and routing inside the customer-controlled environment.
  \item[Model.] Performs bounded language, reasoning, extraction, visual, or
        other computation locally, in customer cloud, or through an optional
        provider.
  \item[Harness.] Prepares tasks, supplies tools, constrains outputs, and
        evaluates behavior as a signed AION capability package.
  \item[Pilot.] Provides conversation, presentation, approval, and control on
        a phone, television, desktop, browser, or authorised device.
  \item[Connectors.] Access authorised files, services, business systems, and
        devices through customer-controlled or explicitly approved boundaries.
  \item[Receipts.] Record proposals, approvals, routes, executions, evidence,
        and outcomes inside the customer's AION brain.
\end{description}

A model adapter may propose content but may not grant itself a capability,
select a private identity, bypass approval, or write unverified statements into
canonical organisational memory.  A capability harness may constrain and use a
model but remains subject to the same identity, policy, evidence, and audit
boundaries as every other AION component.

This arrangement allows the customer to stop using one model or infrastructure
supplier, connect a later Chinese, European, American, or internally developed
model, and retain the accumulated organisational intelligence.

\subsection{Governed remote-compute sequence}

When AION uses a remote model, the complete operating sequence is:

\begin{enumerate}
  \item identify the requesting person, active space, organisation, role, and
        current possession authority;
  \item determine the requested outcome and required capability;
  \item retrieve only information the person and task are authorised to use;
  \item minimise, redact, or replace unnecessary sensitive information;
  \item calculate permitted routes using quality, privacy, residency, latency,
        cost, energy, context, health, and historical performance;
  \item disclose the minimum permitted working context to the selected endpoint;
  \item require a structured response matching the capability contract;
  \item validate structure, claims, citations, constraints, and tool arguments;
  \item convert acceptable output into a governed answer or action proposal;
  \item request possession-bound private approval for consequential work;
  \item execute through a separately authorised connector or device capability;
  \item verify the external result and create an execution receipt;
  \item record bounded outcome evidence for future route selection and learning.
\end{enumerate}

This sequence prevents a model response from being confused with a completed
business action.  A statement such as ``the email was sent'' or ``the event was
booked'' is not accepted until the authorised provider or device returns a
verifiable receipt.

\subsection{Illustrative business operation}

Consider a business owner who asks:

\begin{quote}
Pilot, analyse why sales fell in southern Spain and prepare a recovery plan.
\end{quote}

AION first identifies the owner and the active organisation.  It retrieves only
the authorised sales, marketing, location, customer, and regional evidence.  It
removes information which is unnecessary for the analysis and assesses whether
the local model meets the task's required quality and context threshold.  If the
local route is insufficient, AION may select a qualified model running in the
company's AWS, Google Cloud, NVIDIA, Prem, or other private environment.

The selected model receives a bounded working package and returns a structured
analysis rather than control of the underlying systems.  AION checks the claims
against authorised records, marks unresolved uncertainty, and presents a
verified briefing.  It may prepare proposed actions for Sales, Marketing,
Finance, and Operations, but it executes nothing until the relevant people
approve the exact scopes.  Later outcomes are compared with the plan so that
AION can learn whether the selected model, evidence route, and procedure were
effective.

\subsection{Routing and fallback requirements}

The default routing order is:

\begin{enumerate}
  \item deterministic local capability;
  \item existing AION knowledge, memory, and learned procedures;
  \item a qualified local model;
  \item free or public evidence acquisition;
  \item a customer-configured low-cost provider;
  \item a customer-owned private compute endpoint;
  \item an optional specialist or premium frontier provider;
  \item human review or an honest limitation.
\end{enumerate}

This is a policy order rather than a blind sequence.  A high-sensitivity request
may prohibit all external providers; a long-context but low-sensitivity task may
prefer customer cloud compute; a safety-critical decision may require both
specialist evidence and human review.  A fallback may never cross a privacy,
residency, authority, or spending boundary simply because the preferred model
failed.

\subsection{Deployment acceptance tests}

The deployment system is not complete until it passes the following exercises:

\begin{enumerate}
  \item \textbf{Air-gapped operation:} Pilot, AION memory, the business map, and
        core local capabilities remain useful with the network disabled.
  \item \textbf{No-paid-AI operation:} the system completes representative work
        using local models and permitted public evidence without a premium API.
  \item \textbf{Private-compute operation:} a customer-owned endpoint processes
        a bounded task and returns a contract-valid result without obtaining
        canonical write authority.
  \item \textbf{Premium-provider operation:} an optional provider is used only
        for an eligible request and the disclosure and cost appear in the route
        receipt.
  \item \textbf{Provider-removal operation:} all external provider credentials
        are removed, the AION brain restarts, and identity, memory, business map,
        tasks, policies, evaluations, and receipts remain intact.
  \item \textbf{Model-replacement operation:} a new qualified model is attached
        and continues through the same provider-neutral capability contracts.
  \item \textbf{Migration operation:} the customer moves from a local machine to
        customer cloud or another owned host without creating a second writable
        brain or losing provenance.
  \item \textbf{Boundary operation:} a model-generated instruction attempting to
        grant authority, reveal inaccessible context, or bypass approval fails
        closed and produces an audit event.
\end{enumerate}

\subsection{Strategic conclusion}

Cloud or GPU infrastructure performs selected calculations and returns bounded
results, but AION remains the brain, governor, memory, and operator.  Remote
compute should be treated like a temporary specialist: it receives only the
information required for the authorised assignment, returns work for review,
and does not take ownership of the organisation.

The decisive portability proof is therefore:

\begin{quote}
A customer can remove every model provider, preserve the complete AION brain,
connect a qualified replacement, and continue operating without rebuilding the
organisation's memory, authority, workflows, or business map.
\end{quote}
